% appP_checks.tex -- generated by writing/tools/blueprint_gen.py 1.0 (level 0) on 2026-10-14 from sources 08ca495089; do not edit
\providecommand{\bpNum}[1]{\hyperref[#1]{\ref*{#1}}}%
\providecommand{\bpNumA}[1]{\hyperref[#1]{\textcolor{orange!65!black}{\ref*{#1}}}}%
\providecommand{\bpPage}[1]{\pageref*{#1}}%
\providecommand{\bpLab}[1]{{\ttfamily #1}}%
\providecommand{\bpBrk}{\allowbreak}%
\providecommand{\bpT}[2]{\facttagbox[#1]{\scriptsize #2}}%
\providecommand{\bpU}[1]{{\color{gray!70!black}#1}}%

\section{一致性检查}\label{sec:appP-checks}

\begingroup\footnotesize\setlength{\tabcolsep}{2.5pt}\renewcommand{\arraystretch}{1.06}
\begin{longtable}{@{}>{\raggedright\arraybackslash}p{\dimexpr(\textwidth-6\tabcolsep-0mm)*72/1000\relax}>{\raggedright\arraybackslash}p{\dimexpr(\textwidth-6\tabcolsep-0mm)*257/1000\relax}>{\raggedright\arraybackslash}p{\dimexpr(\textwidth-6\tabcolsep-0mm)*106/1000\relax}>{\raggedright\arraybackslash}p{\dimexpr(\textwidth-6\tabcolsep-0mm)*564/1000\relax}@{}}
\caption{一致性检查的结果（生成程序 blueprint\_gen.py 版本 1.0，日期 2026-10-14，源文件修订 \texttt{08ca495089}，压缩级别 0；E 为错误，W 为警告，I 为信息）}\label{tab:appP-checks}\\
\toprule
\textbf{检查} & \textbf{内容} & \textbf{结果} & \textbf{前三条发现} \\
\midrule
\endfirsthead
\caption[]{一致性检查的结果（续）}\\
\toprule
\textbf{检查} & \textbf{内容} & \textbf{结果} & \textbf{前三条发现} \\
\midrule
\endhead
\midrule
\multicolumn{4}{r@{}}{\footnotesize（续下页）}\\
\endfoot
\bottomrule
\endlastfoot
\mbox{C-P0} & 源文件解析（环境、标题、标签） & W\,3 & \bpLab{\detokenize{lem:}\bpBrk \detokenize{dis-}\bpBrk \detokenize{brackets}}：可选标题的方括号；\newline \bpLab{\detokenize{def:}\bpBrk \detokenize{dis-}\bpBrk \detokenize{latelabel}}：\texttt{\string\label} 不在环境的前三行；\newline ch07\_dis.tex 第 110 行：无标签的 remark 环境 \\
\mbox{C-P1} & 依赖、前提与非微扰地位中的标签均有定义；\texttt{\string\uses} 的格式 & E\,3，W\,1 & \bpLab{\detokenize{lem:}\bpBrk \detokenize{dis-}\bpBrk \detokenize{twouses}}：多于一个 \texttt{\string\uses}；\newline \bpLab{\detokenize{thm:}\bpBrk \detokenize{dis-}\bpBrk \detokenize{nouses}}：缺少 \texttt{\string\uses}；\newline \bpLab{\detokenize{ex:}\bpBrk \detokenize{dis-}\bpBrk \detokenize{seclabel}}：\texttt{\string\uses} 含非陈述标签 \bpLab{\detokenize{sec:}\bpBrk \detokenize{dis-}\bpBrk \detokenize{kin}} \\
\mbox{C-P2} & 规范标签由登记的所属单元写出，且与登记表一致 & E\,7，W\,1，I\,1 & 规范标签 \bpLab{\detokenize{prop:}\bpBrk \detokenize{dis-}\bpBrk \detokenize{missing}}（Ch7）缺失；\newline \bpLab{\detokenize{prop:}\bpBrk \detokenize{dis-}\bpBrk \detokenize{strengthened}}：[A] 与登记表 [P] 不同；\newline \bpLab{\detokenize{prop:}\bpBrk \detokenize{dy-}\bpBrk \detokenize{misplaced}} 由 Ch12 定义，登记的所属单元为 Ch7 \\
\mbox{C-P3} & 无前向依赖 & E\,2，W\,1，I\,1 & \bpLab{\detokenize{lem:}\bpBrk \detokenize{dis-}\bpBrk \detokenize{forward}} 使用后面第~12~章的 \bpLab{\detokenize{thm:}\bpBrk \detokenize{dy-}\bpBrk \detokenize{factorization}}；\newline \bpLab{\detokenize{lem:}\bpBrk \detokenize{appN-}\bpBrk \detokenize{cyc-}\bpBrk \detokenize{a}} 使用同一单元中后出现的 \bpLab{\detokenize{lem:}\bpBrk \detokenize{appN-}\bpBrk \detokenize{cyc-}\bpBrk \detokenize{b}}；\newline \bpLab{\detokenize{lem:}\bpBrk \detokenize{dis-}\bpBrk \detokenize{hidden}} 经 \bpLab{\detokenize{lem:}\bpBrk \detokenize{appN-}\bpBrk \detokenize{hidden}} 间接依赖后面的 \bpLab{\detokenize{thm:}\bpBrk \detokenize{dy-}\bpBrk \detokenize{factorization}} \\
\mbox{C-P4} & 无依赖环 & E\,1 & 依赖环：\bpLab{\detokenize{lem:}\bpBrk \detokenize{appN-}\bpBrk \detokenize{cyc-}\bpBrk \detokenize{a}}、\bpLab{\detokenize{lem:}\bpBrk \detokenize{appN-}\bpBrk \detokenize{cyc-}\bpBrk \detokenize{b}} \\
\mbox{C-P5} & 等级传递：标签不高于前提的内容等级 & E\,3 & \bpLab{\detokenize{lem:}\bpBrk \detokenize{dis-}\bpBrk \detokenize{forward}} [R] 使用 \bpLab{\detokenize{thm:}\bpBrk \detokenize{dis-}\bpBrk \detokenize{ope}} [A]；\newline \bpLab{\detokenize{lem:}\bpBrk \detokenize{dis-}\bpBrk \detokenize{forward}} [R] 使用 \bpLab{\detokenize{thm:}\bpBrk \detokenize{dy-}\bpBrk \detokenize{factorization}} [A]；\newline \bpLab{\detokenize{prop:}\bpBrk \detokenize{dy-}\bpBrk \detokenize{ctech}} [A] 使用假设 \bpLab{\detokenize{asm:}\bpBrk \detokenize{appN-}\bpBrk \detokenize{brst-}\bpBrk \detokenize{nonlocal}} [C] \\
\mbox{C-P6} & 正面陈述不依赖 [V] 陈述 & E\,1 & \bpLab{\detokenize{thm:}\bpBrk \detokenize{dis-}\bpBrk \detokenize{usesv}} 使用 [V] 陈述 \bpLab{\detokenize{prop:}\bpBrk \detokenize{dis-}\bpBrk \detokenize{violated}} \\
\mbox{C-P7} & [R] 的限定语逐字一致且彼此相容 & E\,2，I\,1 & \bpLab{\detokenize{prop:}\bpBrk \detokenize{dis-}\bpBrk \detokenize{rdef}} [R, def] 使用 \bpLab{\detokenize{lem:}\bpBrk \detokenize{dis-}\bpBrk \detokenize{ptlemma}} [R, PT]；\newline \bpLab{\detokenize{prop:}\bpBrk \detokenize{dis-}\bpBrk \detokenize{notverbatim}}：限定语句不合规定 \\
\mbox{C-P8} & $R_{\mathrm{cond}}$ 的前提句完整 & E\,4 & \bpLab{\detokenize{thm:}\bpBrk \detokenize{appN-}\bpBrk \detokenize{structure}}：假设 \bpLab{\detokenize{asm:}\bpBrk \detokenize{appN-}\bpBrk \detokenize{wrongtag}} 不在前提句中；\newline \bpLab{\detokenize{thm:}\bpBrk \detokenize{appN-}\bpBrk \detokenize{structure}}：前提 \bpLab{\detokenize{lem:}\bpBrk \detokenize{appN-}\bpBrk \detokenize{cyc-}\bpBrk \detokenize{a}} 不在 \texttt{\string\uses} 中；\newline \bpLab{\detokenize{thm:}\bpBrk \detokenize{appN-}\bpBrk \detokenize{noprem}}：缺少前提句 \\
\mbox{C-P9} & 每条编号假设至少被使用一次 & E\,1 & 假设 \bpLab{\detokenize{asm:}\bpBrk \detokenize{appN-}\bpBrk \detokenize{unused}} 未被任何陈述使用 \\
\mbox{C-P10} & 「被使用于」与计算的直接使用者一致 & W\,4 & \bpLab{\detokenize{asm:}\bpBrk \detokenize{appN-}\bpBrk \detokenize{collins}}：登记表缺 4、多 0；\newline \bpLab{\detokenize{asm:}\bpBrk \detokenize{appN-}\bpBrk \detokenize{collins}}：正文缺 5、多 0；\newline \bpLab{\detokenize{asm:}\bpBrk \detokenize{appN-}\bpBrk \detokenize{values}}：登记表缺 1、多 0 \\
\mbox{C-P11} & 主表各行的标签存在且严格性标签一致 & E\,2 & 主表 III.4：\bpLab{\detokenize{thm:}\bpBrk \detokenize{dis-}\bpBrk \detokenize{nonexistent}} 未定义；\newline 主表 II.2：\bpLab{\detokenize{thm:}\bpBrk \detokenize{dis-}\bpBrk \detokenize{ope}} 的标签与正文不同 \\
\mbox{C-P12} & 无重复标签 & E\,2，W\,1 & \bpLab{\detokenize{fig:}\bpBrk \detokenize{synth-}\bpBrk \detokenize{same}} 定义了 2 次；\newline \bpLab{\detokenize{sec:}\bpBrk \detokenize{appP-}\bpBrk \detokenize{register}} 由生成程序定义，单元中不应再定义；\newline \bpLab{\detokenize{sec:}\bpBrk \detokenize{appP-}\bpBrk \detokenize{register}} 已由单元定义，生成的节标题省略 \\
\mbox{C-P13} & 无待核实标记与已撤回的标签 & W\,2 & ch12\_dy.tex 第 79 行：待核实注释；\newline ch12\_dy.tex 第 80 行：引用已撤回的标签 \bpLab{\detokenize{thm:}\bpBrk \detokenize{appN-}\bpBrk \detokenize{oneloop}} \\
\mbox{C-P14} & 主表“瓶颈”列包含计算的瓶颈 & W\,1 & 主表 III.2（\bpLab{\detokenize{prop:}\bpBrk \detokenize{dy-}\bpBrk \detokenize{model}}）：缺 1、多 1 \\
\mbox{C-P15} & 强子层面陈述的非微扰地位行 & E\,2 & \bpLab{\detokenize{prop:}\bpBrk \detokenize{dy-}\bpBrk \detokenize{npline}}：非微扰地位行含非全局标签 \bpLab{\detokenize{asm:}\bpBrk \detokenize{appN-}\bpBrk \detokenize{collins}}；\newline \bpLab{\detokenize{prop:}\bpBrk \detokenize{dy-}\bpBrk \detokenize{hadron}}：强子层面陈述缺非微扰地位行 \\
\mbox{C-P16} & 「地位与缺口」段落存在并点名瓶颈 & E\,1，W\,2 & \bpLab{\detokenize{conj:}\bpBrk \detokenize{dy-}\bpBrk \detokenize{coloured}} [C]：缺「地位与缺口」；\newline \bpLab{\detokenize{prop:}\bpBrk \detokenize{method-}\bpBrk \detokenize{gapless}}：「地位与缺口」未点名瓶颈（\bpLab{\detokenize{lem:}\bpBrk \detokenize{method-}\bpBrk \detokenize{ward-}\bpBrk \detokenize{multi}} 等）；\newline \bpLab{\detokenize{prop:}\bpBrk \detokenize{dy-}\bpBrk \detokenize{late}}：「地位与缺口」未紧随环境 \\
\mbox{C-P17} & 具名缺口被点名 & W\,1 & \bpLab{\detokenize{prop:}\bpBrk \detokenize{method-}\bpBrk \detokenize{gapless}}：未点名具名缺口 \bpLab{\detokenize{asm:}\bpBrk \detokenize{appN-}\bpBrk \detokenize{brst-}\bpBrk \detokenize{nonlocal}} \\
\end{longtable}
\endgroup
完整的发现清单（按单元分组）见生成程序的报告文件 \texttt{appP\_report.txt}。
\FloatBarrier\label{appP:end}
